\section{On the Equivalence of \LCC and \CP}
\label{sec:equivalence}

In this section we sketch in a more detailed way the relationship between \LCC
and \CP. More specifically, we show that there is an isomorphism between
(well-typed) \CP processes and (well-typed) \LCC processes, up to renaming of
bound names. Such tight correspondence is essentially the same studied by
Kobayashi~\cite{Kobayashi02b} and Dardha et al.~\cite{DardhaGiachinoSangiorgi17}
between session calculi and the linear $\pi$-calculus and by Sano et
al.~\cite{SanoKavanaghPientka23} between \CP and \SCP.
%
For the sake of illustration, we only focus on the process forms corresponding
to the multiplicative connectives $\Ten$ and $\Par$, leaving the extension to
the additive connectives as a straightforward exercise.

If we let $\ClassicalP$ and $\ClassicalQ$ (in sans-serif font) range over \CP
processes and write $\wtcp\ClassicalP\Context$ for a \CP typing judgment, in \CP
the process forms representing input/output operations on sessions are
$\Join\x\y{}.\ClassicalP$ and $\Fork\x\y{}\ClassicalP\ClassicalQ$ and their
typing rules are formulated thus:
\[
    \inferrule[\StandardForkRule]{
        \wtcp\ClassicalP{\ContextC, y : A}
        \\
        \wtcp\ClassicalQ{\ContextD, x : B}
    }{
        \wtcp{\Fork\x\y{}\ClassicalP\ClassicalQ}{\ContextC, \ContextD, x : A \Ten B}
    }
    \qquad
    \inferrule[\StandardJoinRule]{
        \wtcp\ClassicalP{\Context, y : A, x : B}
    }{
        \wtcp{\Join\x\y{}.\ClassicalP}{\Context, x : A \Par B}
    }
\]

The key difference between these forms and the corresponding ones in \LCC is the
\emph{absence} of an explicit continuation channel. Instead, the session
endpoint $x$ is being used by the continuation processes $Q$ and $P$ to
implement the rest of the session.

\newcommand\encode[1]{\varepsilon(#1)}
\newcommand\decode[1]{\delta(#1)}

Every (well-typed) \CP process can be \emph{encoded} into a (well-typed) \LCC
process by introducing an explicit continuation channel that does not clash with
other channel names. Since we identify processes up to bound names, we may pick
as the name of the continuation channel the very same name of the session that
it continues and express the encoding as a function $\encode\cdot$ that includes
the following equations:
\[
    \encode{\Fork\x\y{}\ClassicalP\ClassicalQ} = \Fork\x\y\x{\encode\ClassicalP}{\encode\ClassicalQ}
    \qquad
    \encode{\Join\x\y{}.\ClassicalP} = \Join\x\y\x.\encode\ClassicalP
\]

Every (well-typed) \LCC process can be \emph{decoded} into a (well-typed) \CP
process by renaming each continuation channel so that it matches the name of the
session it continues:
\[
    \decode{\Fork\x\y\z{P}{Q}} = \Fork\x\y{}{\decode{P}}{\decode{Q\subst\z\x}}
    \qquad
    \decode{\Join\x\y\z.R} = \Join\x\y{}.\decode{R\subst\z\x}
\]

It is easy to establish that $\encode\cdot$ and $\decode\cdot$ are one the
inverse of the other and that they preserve typing, in the sense that:
\begin{itemize}
\item If $\wtcp\ClassicalP\Context$, then $\wtp{\encode\ClassicalP}\Context$;
\item If $\wtp{P}\Context$, then $\wtcp{\decode{P}}\Context$.
\end{itemize}
